% Pursuing Stacks (cotes 134-2 to 134-8): archive cote 134-7, site page 24 (PDF page 25), typescript page 466.
% typescript. Draft compared: ps6.tex:1106-1176.
% AI-transcribed by Claude Opus 5.5 (claude-opus-5-5, adaptive thinking, effort high), Message Batches API, batch ps3 (msgbatch_01XdFmuVVBAcgc7EFxdtCqsj), 2026-11-12; not yet checked by eye.
\scanpage{134-7/24}
\mspage{466}

27.8.

\margin{113.} What I was thinking of last night (see last sentence) is that
whereas for total homology (not for the separate H$_i$'s) we have the com-%
prehensive formula
\[ \text{A)} \qquad \mathrm{LH}_{.}(X_{*}\otimes_k k')\hand{\simeq}\mathrm{LH}_{.}(X_{*}) \overset{\hand{L}}{\otimes}_k k' \quad \hand{\unsure{(}}\text{ in } D_{.}(k')) , \]
\hand{(}where $\overset{\hand{L}}{\otimes}_k$ denotes the left derive\hand{d} functor of the ring extension functor
\inserted{for} $k\to k'$, and \inserted{$X'_{*}=$} $X_{*}\otimes_k k'$ denotes ring extension for the semisimplicial unipo-%
tent bundle $X_{*}$\hand{,)} for homotopy modules we should have the term-by-term
isomorphisms
\[ \text{B)} \qquad \hand{\pi}_i(X_{*}\otimes_k k')\hand{\unsure{\leftarrow}}\hand{\pi}_i(X_{*})\otimes_k k' \quad . \]
\margin{\hand{$\leftarrow$} (circled)}
This however was pretty rash indeed (it was time to go to sleep I guess !\hand{)}
\hand{\unsure{.}} \over{w}{W}hereas the map on sections
\[ X_{*}(k)\longrightarrow X'_{*}(k') \]
does induce a map
\[ \text{B')} \qquad \hand{\pi}_i(X_{*})\deleted{\illegible{1}}\hand{\simeq}\hand{\pi}_i(X_{*}(k)) \longrightarrow \hand{\pi}_i(X'_{*})\hand{\simeq}\hand{\pi}_i(X'_{*}(k')) \]
which surely is k-linear, and \inserted{hence} induces a map B), this map is certainly
\uline{not} an isomorphism without some flatness restriction either on k' over k,
or on the k-modules $\hand{\pi}_j(X_{*})$ for $\hand{j<i}$, as we had noted already three days
ago when looking at the case when $X_{*}$ comes from a chain complex $M_{\hand{\bullet}}$ in
$(\mathrm{Ab}_k)$ (with projective coefficients say), and hence $X'_{*}$ comes from
\[ M_{\over{*}{\bullet}}' = M_{.}\otimes_k k' \quad . \]
If we look at the description of the k-\uline{module} $\hand{\pi}_i(X_{*}) = \hand{\pi}_i$ in terms of a
Postnikov dévissage of $X_{*}$ , we should recall that the \xout{group object}
semisimplicial gro\xout{y}up object $M(i)_{*}$ which inters into the picture as the
i.th step fiber is \uline{not} the one defined \uline{directly} (via the Kan-Dold-Puppe
functor) by $\hand{\pi}_i$ placed in degree i , but \inserted{rather} by the chain complex \xout{with}
$M(i)_{\over{*}{\bullet}}$ \uline{with projective components}, obtained by taking \inserted{first} a shifted projec-%
tive resolution of $\hand{\pi}_i$ . Thus, by ring extension we get from this dévis-%
sage of $X_{*}$ another one of $X'_{*}$ , whose successive fibers are
\[ M(i)'_{*} = M(i)_{*}\otimes_k k' \quad , \]
\xout{with canonic augmentations}
corresponding to the chain complexes
\[ M(i)_{.}' = M(i)_{.}\otimes_k k' \quad . \]

\margin{24}
